\clearpage\fancyhead[L]{\small REVIEW.}\pdfbookmark[0]{Reviews. 1815 July 1.}{anzeige-1815-pfaff}\SourceLeaf{240}
\begin{center}\textbf{R E V I E W S.}\par\bigskip
\rule{.85\linewidth}{.4pt}\par
Göttingische gelehrte Anzeigen.\quad1815 July 1.\par
\rule{.85\linewidth}{.4pt}\end{center}
Professor \textsc{Gauss} has presented to the Royal Society a manuscript paper by Professor \textsc{Pfaff} of Halle, entitled:
\begin{center}\textit{Methodus generalis, aequationes differentiarum partialium, nec non aequationes differentiales vulgares, utrasque primi ordinis, inter quotcunque variabiles, complete integrandi.}\end{center}
The theory of first-order partial differential equations is known to owe two important extensions to \textsc{Lagrange}: their general integration either when, with any number of variables, they contain the partial derivatives only linearly, or, without restriction on their form, when there are only three variables. Until now, work had not really advanced beyond these two cases, and the method used by \textsc{Lagrange} for partial differential equations in three variables would present particular difficulties with a larger number. We must therefore regard it as a remarkable enrichment of the integral calculus that the perceptive author of the present paper has succeeded in finding the general integration of first-order partial differential equations for \textit{every} number of variables. In this investigation he has chosen a path of his own\SourceLeaf{241}%
 and connected it with another, no less interesting branch of the integral calculus: the theory of ordinary differential equations (of the first order) between more than two variables, whose true nature, as is known, \textsc{Monge} was the first to teach us, although this geometer completed their integration only for the simplest cases. Mr. \textsc{Pfaff} gives the general integration of this kind of differential equation, and that of partial differential equations then appears only as a special case of it. These investigations \textit{presuppose} the general integration of differential equations of every order between two variables. This is entirely appropriate, just as the general solution of algebraic equations is postulated in the higher parts of mathematics.

We believe we shall do those interested in higher mathematics a welcome service by placing this beautiful extension of the integral calculus at their disposal through the present review. An extract from the paper, which occupies 144 quarto pages, in which we followed the author step by step and passed over nothing essential, would, of course, far exceed the space allowed us. We shall therefore try, keeping only to the \textit{subject itself}, to bring out the essentials in a somewhat different presentation, so that knowledgeable readers can fully make them their own.

The main operation in the whole undertaking must be regarded as the reduction of a differential expression
\[
\Omega=p\,\mathrm dx+p'\mathrm dx'+p''\mathrm dx''+\text{etc.}+p^{(n-1)}\mathrm dx^{(n-1)}
\]
where each coefficient $p,p',p''$ etc. is a function of the variables $x,x',x''$ etc., to the form
\[
\Omega=\lambda(q\,\mathrm dy+q'\mathrm dy'+q''\mathrm dy''+\text{etc.}+q^{(n-2)}\mathrm dy^{(n-2)});
\]
such that $\lambda,y,y',y''$ etc. are functions of $x,x',x''$ etc., whereas $q,q',q''$ etc. are functions of $y,y',y''$ etc., and the number of these latter variables is one less than the number of variables $x,x',x''$ etc.

This transformation, which we shall denote by (I) for brevity, is restricted to the case where $n$ is an even number. We shall explain below\SourceLeaf{242}%
 how it is to be carried out, and shall assume it here for the moment so as not to interrupt the overview.
\par
In the case where $n$ is an odd number, transformation I would be possible only under special conditions between the coefficients $p,p',p''.\,.\,.$; in general, however, $\Omega$ can in this case be brought to the form
\[
p^{*}\mathrm{d}x+\lambda(q\mathrm{d}y+q'\mathrm{d}y'+q''\mathrm{d}y''+\mathrm{etc.}
+q^{(n-3)}\mathrm{d}y^{(n-3)})
\]
For the moment, regard $x$ in $\Omega$ as constant, and under this assumption use (I) to transform
\[
p'\mathrm{d}x'+p''\mathrm{d}x''+p'''\mathrm{d}x'''+\mathrm{etc.}
+p^{(n-1)}\mathrm{d}x^{(n-1)}
\]
where the number of variables $x',x'',x'''.\,.\,.$ will now be even, into
\[
\lambda(q\mathrm{d}y+q'\mathrm{d}y'+q''\mathrm{d}y''+\mathrm{etc.}
+q^{(n-3)}\mathrm{d}y^{(n-3)})=\lambda\Omega'
\]
Here, then, $q,q',q''$ etc. will be functions of $y,y',y''$ etc., while these, like $\lambda$, will be functions of $x,x',x'',x'''.\,.\,.\,.$. Of these quantities, however, the first, $x$, must be treated as constant in order to obtain, by expanding $\lambda\Omega'$,
\[
p'\mathrm{d}x'+p''\mathrm{d}x''+p'''\mathrm{d}x'''+\mathrm{etc.}
\]
The additional term that arises if $x$ is also regarded as variable in that expansion is
\[
=(q\cdot\frac{\mathrm{d}y}{\mathrm{d}x}
+q'\cdot\frac{\mathrm{d}y'}{\mathrm{d}x}
+q''\cdot\frac{\mathrm{d}y''}{\mathrm{d}x}+\mathrm{etc.})\lambda\mathrm{d}x
\]
Thus, to obtain the form above, one need only set
\[
p^{*}=p-\lambda(q\cdot\frac{\mathrm{d}y}{\mathrm{d}x}
+q'\cdot\frac{\mathrm{d}y'}{\mathrm{d}x}
+q''\cdot\frac{\mathrm{d}y''}{\mathrm{d}x}+\mathrm{etc.})
\]
We shall denote this transformation of $\Omega$ into $p^{*}\mathrm{d}x+\lambda\Omega'$, which is restricted to odd values of $n$, by II. The same reduction can evidently be applied again to $\Omega'$, putting
\[
\begin{aligned}
&\Omega'=q^{*}\mathrm{d}y+\lambda'(r\mathrm{d}z+r'\mathrm{d}z'+r''\mathrm{d}z''+\mathrm{etc.}\\&\qquad+r^{(n-5)}\mathrm{d}z^{(n-5)})\\
&=q^{*}\mathrm{d}y+\lambda'\Omega''
\end{aligned}
\]
\SourceLeaf{243}
and then again $\Omega''=r^{*}\mathrm{d}z+\lambda''\Omega'''$, until one finally reaches an expression containing only one variable. Thus $\Omega$ has been brought to the form
\[
p^{*}\mathrm{d}x+\lambda q^{*}\mathrm{d}y+\lambda\lambda'r^{*}\mathrm{d}z+\mathrm{etc.}
\]
or to the form
\[
P\mathrm{d}x+Q\mathrm{d}y+R\mathrm{d}z+\mathrm{etc.}
\]
where the number of variables $x,y,z$ etc. $=\frac12(n+1)$, and all $n$ quantities $y,z$ etc., $P,Q,R$ etc. will be functions of $x,x',x''$ etc. Let this reduction procedure be denoted by III.

If, in the case where an even number of variables was originally given, this procedure III is applied to the expression
\[
q\mathrm{d}y+q'\mathrm{d}y'+q''\mathrm{d}y''+\mathrm{etc.}
\]
obtained by reduction I, then
\[
\begin{aligned}&\Omega=p\mathrm{d}x+p'\mathrm{d}x'+p''\mathrm{d}x''+\mathrm{etc.}
+p^{(n-1)}\mathrm{d}x^{(n-1)}\\&\qquad\text{into the form}\end{aligned}
\]
\[
Q\mathrm{d}y+R\mathrm{d}z+\mathrm{etc.}
\]
so that the number of variables $y,z$ etc. $=\frac12 n$, and all $n$ quantities $Q,R$ etc., $y,z$ etc. become functions of $x,x',x''$ etc. Let this reduction be denoted by IV.

This general transformability of differential expressions according to III and IV is a theorem as new as it is remarkable. Although not explicitly stated in Mr. \textsc{Pfaff}'s paper, it follows readily from the investigations there.

The solutions of the problems mentioned at the beginning of this review can now easily be derived from it.

\Needspace{17\baselineskip}
\medskip
1) To integrate the differential equation
\[
0=p\mathrm{d}x+p'\mathrm{d}x'+p''\mathrm{d}x''+\mathrm{etc.}
\]
or
\[
0=\Omega
\]
\SourceLeaf{244}
where $p,p',p''$ etc. are given functions of the $n$ variables $x,x',x''$ etc., one will, if $n$ is even, use IV to put
\[
\Omega=Q\mathrm{d}y+R\mathrm{d}z+S\mathrm{d}u+\mathrm{etc.}
\]
where $Q,y.R,z,S,u$ etc. will together be $n$ given functions of $x,x',x''$. The differential equation
\[
0=Q\mathrm{d}y+R\mathrm{d}z+S\mathrm{d}u+\mathrm{etc.}
\]
will thus be equivalent to the given one, and its most general integration will be contained in the following system of $\frac12 n$ equations:
\[
\begin{aligned}&0=\varphi(y,z,u\ \mathrm{etc.}),\\&
\frac1Q\cdot\frac{\mathrm{d}\varphi(y,z,u.\,.\,.)}{\mathrm{d}y}
\\&=\frac1R\cdot\frac{\mathrm{d}\varphi(y,z,u.\,.\,.)}{\mathrm{d}z}
\\&=\frac1S\cdot\frac{\mathrm{d}\varphi(y,z,u.\,.\,.)}{\mathrm{d}u}=\mathrm{etc.}\end{aligned}
\]
where $\varphi$ represents an arbitrary function, and the derivatives are, of course, \textit{partial}. Insofar as the equation
$0=\varphi(y,z,u\ \mathrm{etc.})$
allows $y$ to be determined in terms of the others, the solution can also be represented by the following equations:
\[
\begin{gathered}
y=\psi(z,u.\,.\,.\,.)\\
\frac{\mathrm{d}\psi(z,u.\,.\,.\,.)}{\mathrm{d}z}=-\frac RQ\\
\frac{\mathrm{d}\psi(z,u.\,.\,.\,.)}{\mathrm{d}u}=-\frac SQ\\
\text{etc.}
\end{gathered}
\]
Strictly speaking, however, this solution would be less general, since the arbitrary function $\varphi(y,z,u.\,.\,.\,.)$ also includes functions in which $y$ does not occur.

\Needspace{17\baselineskip}
\medskip
2) To integrate the same differential equation in the case where $n$ is odd, one will put $\Omega$ into the following form according to III:
\[
\Omega=P\mathrm{d}x+Q\mathrm{d}y+R\mathrm{d}z+\mathrm{etc.}
\]
where $P,Q,y,R,z$ etc. will together be $n$ given functions of $x,x',x''$ etc. The most general integration of the differential equation $\Omega=0$ then rests on the following system of $\frac12(n+1)$ equations:
\[
\begin{aligned}&0=\varphi(x,y,z.\,.\,.),\\&
\frac1P\cdot\frac{\mathrm{d}\varphi(x,y,z.\,.\,.)}{\mathrm{d}x}
=\frac1Q\cdot\frac{\mathrm{d}\varphi(x,y,z.\,.\,.)}{\mathrm{d}y}
\\&=\frac1R\cdot\frac{\mathrm{d}\varphi(x,y,z.\,.\,.)}{\mathrm{d}z}=\mathrm{etc.}\end{aligned}
\]
\par\SourceLeaf{245}
3) The general integration of a given first-order partial differential equation, that is, of a finite equation between the partial derivatives
\[
\frac{\mathrm{d}x}{\mathrm{d}x'}=p',\qquad
\frac{\mathrm{d}x}{\mathrm{d}x''}=p'',\qquad
\frac{\mathrm{d}x}{\mathrm{d}x'''}=p''',\quad\text{etc.}
\]
and $x,x',x'',x'''$ etc. (where $x$ represents a function of the $m$ variables $x',x'',x'''$ etc. that is still to be determined), is nothing other than the general integration of the ordinary differential equation
\[
0=-\mathrm{d}x+p'\mathrm{d}x'+p''\mathrm{d}x''+p'''\mathrm{d}x'''+\text{etc.}
\]
For, since that finite equation allows one of the quantities $p',p'',p'''$ etc., for example $p'$, to be represented as a function of the others, $p'',p'''$ etc., and of $x,x',x'',x'''$ etc., the differential equation just given is to be regarded as one between the $2m$ variables $x,x',x''$ etc., $p'',p'''$ etc., in which the differentials $\mathrm{d}p'',\mathrm{d}p'''$ etc. have the coefficient $0$. Thus, to carry out the integration, one will bring the differential expression
\[
-\mathrm{d}x+p'\mathrm{d}x'+p''\mathrm{d}x''+p'''\mathrm{d}x'''+\text{etc.}
\]
to the form
\[
Q\mathrm{d}y+R\mathrm{d}z+S\mathrm{d}u+\text{etc.}
\]
where the $2m$ quantities $Q,y,R,z,S.\ u$ etc. will be known functions of $x,x',x''\, .\, .\, .\, .\, p'',p'''\, .\, .\, .\, .$. The integration is then contained in the same system of equations as above (1), and if one imagines $p'',p'''$ etc. eliminated from them, one finite equation between $x,x',x''.x'''$ etc. remains. The \emph{actual} elimination can, of course, be carried out only insofar as \emph{determinate} functions are taken for $\varphi$. This circumstance, however, rests on the nature of the problem and not on the imperfection of analysis, which, as long as it remains general, can give the solution only in that form.

It is, moreover, immediately apparent that in a similar way the integration of \emph{several} simultaneous partial differential equations is within our power.


